\section{结论}
\label{sec:conclusion}

本文的主要实验结果均应首先理解为选择器解析层面的诊断结果，而非工业部署中的端到端通过率。在缺乏行为验证器（oracle）的部署条件下，LLM 驱动的 Playwright 自愈与“测试通过但真实回归被遮蔽”之间仅一步之遥。本文在对抗性极限下测量到 \texttt{qwen2.5-coder:7b} 的 78.7\% 隐性误修复上界，且提示级放弃指令逐例无效；\texttt{gpt-4o} 与 \texttt{deepseek-v4-flash} 分别能将误修复从 84.0\% 降至 68.0\%、从 80.0\% 降至 70.7\%，但均留下多数错误可解析补丁。HealReact 的 L1/L3 静态基底使 koenig 数据集中 80.6\% 的真实失效在静态层面可达；一个小型本地修复器在可达用例上的 top-1 选择器精确匹配率达到 77.3\%；两个低开销确定性手段给出了统计未显著、且构成依赖底层模型的局部改进——强锚点后置过滤在 Qwen 与 DeepSeek 上均零回归，testId 加权则未在 DeepSeek 上复现净增益。三项结果合起来，为后续工作中评估行为重放验证器给出了直接的经验性动机，也给出了一条由锚点文化差异所限制的、可证伪的跨应用泛化假设。
